\subsection{Promedio de puntos}

Los parámetros \(k_p\) y \(T_d\) se determinan a partir de la gráfica de la
respuesta del sistema. El valor \(A k_p\) corresponde al nivel al que tiende
la salida, mientras que \(T_d\) representa el tiempo que el sistema tarda en
responder tras la activación de la entrada en el instante \(b\).

La constante de tiempo se calcula considerando el intervalo posterior a
\(T_d + b\) necesario para que el sistema alcance porcentajes clave del
valor final, como \qty{28}{\percent}, \qty{50}{\percent} (vida media),
\qty{63}{\percent}, \qty{86}{\percent} y \qty{95}{\percent}. Los tiempos
obtenidos para cada porcentaje se promedian, proporcionando una aproximación
de la constante característica del sistema.

\subsection{Método de dos puntos (Alfaro \cite{ruiz-2011})}

Este método emplea dos tiempos específicos, \(t_1\) y \(t_2\), que corresponden
a los instantes en que el sistema alcanza los porcentajes de \qty{25}{\percent}
y \qty{75}{\percent} del valor final, respectivamente. La ganancia \(k_p\) se
determina de manera gráfica, mientras que las constantes \(T_p\) y \(T_d\) se
calculan mediante las expresiones dadas en las
\cref{eq:Tp-Alfaro,eq:Td-Alfaro}.

\begin{gather}
  T_p = -0.910 t_1 + 0.910 t_2 \label{eq:Tp-Alfaro} \\
  T_d = 1.262 t_1 - 0.262 t_2 \label{eq:Td-Alfaro}
\end{gather}

\subsection{Método de Stark}

\subsection{Método de Jahanmiri y Fallahi \cite{jahanmiri-1997}}
